\begin{python0}
from solutions import *; clear()
\end{python0}
\sectionthree{Default value parameter and nonconstant value}

One more example. Recall we have the \texttt{isprime()} that returns \texttt{true} if the values passed in is a prime number, otherwise \texttt{false} is returned.

\begin{consolethree}[escapeinside=||]
bool isprime(int n)
{
    for (int i = 2; i <= sqrt((double)n); i++)
    {
        if (n % i == 0) return false;
    }
    return true;
}
\end{consolethree}

Suppose for some reason you know something about n: if there is a prime divisor of n, then it must be at least 11. You might want to tell \texttt{isprime()} to start the search for a divisor from 11.

Let's modify the function preserving the older use:

\begin{consolethree}[escapeinside=||]
#include <iostream>
#include <cmath>

bool isprime(int n, int start = 2)
{
    for (int i = start; i <= sqrt((double)n); i++)
    {
        if (n % i == 0) return false;
    }
    return true;
}

int main()
{
    std::cout << isprime(5) << '\n'
              << isprime(5, 3) << '\n'
              << isprime(169, 11) << std::endl;
    return 0;
}
\end{consolethree}

Now let's try to generalize the upper bound of the search. Let's try this:

\begin{consolethree}[escapeinside=||]
#include <iostream>
#include <cmath>

bool isprime(int n,
             int start = 2,
             |\textbf{int end = sqrt((double)n))}|
{
    for (int i = start; i <= end; i++)
    {
        if (n % i == 0) return false;
    }
    return true;
}

int main()
{
    std::cout << isprime(5) << '\n'
              << isprime(5, 3) << '\n'
              << isprime(169, 11) << std::endl;
    return 0;
}
\end{consolethree}

Does it work?

It \EMPHASIZE{won't work} because remember you can only use constants for default values. So how do you set a parameter to ``default value'' that depends on an expression like the above? This is how you do it:

\begin{consolethree}[escapeinside=||]
#include <iostream>
#include <cmath>

bool isprime(int n, int start = 2, int end = -1)
{
    |\textbf{if (end == -1) end = sqrt(double(n));}|
    for (int i = start; i <= end; i++)
    {
        if (n % i == 0) return false;
    }
    return true;
}

int main()
{
    std::cout << isprime(5) << '\n'
              << isprime(5, 3) << '\n'
              << isprime(169, 11) << '\n'
              << isprime(21, 3, 7);
    return 0;
}
\end{consolethree}

Of course in this case we must choose a reasonable default value for parameter \texttt{end}. In the above case, I choose a value the user won't be using. This is similar to the use of the sentinel values in terminating a while-loop.

Note that a default values can be from a function call as long as the function call involves only constants. Make sure you run this:

\begin{consolethree}[escapeinside=||]
int g(int x)
{
    return x * x;
}

void f(int x, int y = g(42))
{}
\end{consolethree}


